% CS F407 Artificial Intelligence: final project report template.
% Limit: 2 pages, including everything except references.
% Do not change the font size, margins or spacing.
\documentclass[11pt,a4paper,twocolumn]{article}

\usepackage[margin=2cm,columnsep=0.8cm]{geometry}
\usepackage[T1]{fontenc}
\usepackage{lmodern}
\usepackage{microtype}
\usepackage{amsmath,amssymb}
\usepackage{graphicx}
\usepackage{booktabs}
\usepackage{tabularx}
\usepackage{enumitem}
\usepackage{titlesec}
\usepackage[hidelinks]{hyperref}

\titleformat{\section}{\normalsize\bfseries}{\thesection}{0.6em}{}
\titlespacing*{\section}{0pt}{1.2ex plus 0.3ex}{0.6ex}
\setlist{nosep,leftmargin=*}
\setlength{\parskip}{0.3ex}

\newcommand{\coursename}{CS F407 Artificial Intelligence, AY 2026--27 Semester I}
\newcommand{\teamno}{Team XX}
\newcommand{\githuburl}{https://github.com/username/repository}

\begin{document}

\twocolumn[{%
\begin{center}
  {\large\bfseries Project Title: Concise and Informative}\\[1ex]
  \teamno\\
  Member One (2022A7PS0001G), Member Two (2022A7PS0002G), Member Three (2022A7PS0003G)\\[0.5ex]
  {\small \coursename, BITS Pilani, K.K. Birla Goa Campus}\\[1.5ex]
  \begin{minipage}{0.85\textwidth}
    \small
    \centerline{\bfseries Abstract}\vspace{0.5ex}
    In 100--150 words: the problem, why it matters, what you did, the main result
    (with a number, if you have one), and the main conclusion.
  \end{minipage}
\end{center}
\vspace{1.5ex}
}]

\section{Introduction}
State the problem and its motivation. Say clearly what the input and output of your system
are, and list your contributions in one or two sentences.

\section{Background and related work}
Briefly describe the AI techniques you build on (for example, search, SAT, game playing,
planning, probabilistic models, LLMs) and the closest existing work, with
citations~\cite{russell2021aima}.

\section{Method}
Describe your approach in enough detail for someone else to reproduce it. Use equations,
pseudocode or a small figure where it helps. If an LLM is part of the method, give the
model, how it is used (prompting, tool use, fine-tuning), and any prompts that matter.

\section{Experimental setup}
Datasets or problem instances, baselines, evaluation metrics, and the hardware and software
used.

\section{Results}
Report your main results in a table or a figure and explain what they show.

\begin{table}[h]
\centering
\small
\begin{tabular}{lcc}
\toprule
Method & Metric 1 & Metric 2 \\
\midrule
Baseline & 0.00 & 0.00 \\
Ours & 0.00 & 0.00 \\
\bottomrule
\end{tabular}
\caption{Main results. Say what the metrics are and whether higher or lower is better.}
\label{tab:results}
\end{table}

\section{Discussion and limitations}
Interpret the results. Where does your approach fail, and why? What would you do next?

\section{Conclusion}
Two or three sentences summarising what you found.

\section*{Individual member contributions}
{\small
\begin{tabularx}{\linewidth}{@{}l X@{}}
\toprule
Name & Contribution \\
\midrule
Member One & For example: problem formulation, search implementation, Sections 1 and 3. \\
Member Two & For example: LLM pipeline, experiments, Section 5. \\
Member Three & For example: data preparation, evaluation, Sections 4 and 6. \\
\bottomrule
\end{tabularx}}

\smallskip
\noindent All members have read this report and agree with the statement above.

\section*{Generative AI disclosure}
State which generative AI tools you used (name and version), what you used them for
(for example, writing or debugging code, editing text, generating data), and which parts of
the report or code they affected. State that you checked all such output and take
responsibility for it. If you used none, write: ``No generative AI tools were used in this
project.'' Use of an LLM as part of the method itself goes in Section~3, not here.

\section*{Code availability}
\noindent Code and data:\\
\url{\githuburl}\\
The repository must be public and must contain a README with instructions to reproduce the
results reported here.

\begin{thebibliography}{9}
\small
\bibitem{russell2021aima}
S.~Russell and P.~Norvig.
\newblock \emph{Artificial Intelligence: A Modern Approach}, 4th edition.
\newblock Pearson, 2021.
\end{thebibliography}

\end{document}
